\documentclass{article}
\usepackage{amsmath,amssymb}
\begin{document}
\section*{A symmetry obstruction for $N=27$, $t=6$}
This is a partial exclusion. It does not prove nonexistence of an arbitrary
sixth-order anticoherent state. The universal sixth-order exclusion, and hence
$t_{\max}(27)=5$, is proved separately by the six-copy certificate in
\texttt{../n27-t6/}.

\paragraph{Claim.}
A pure symmetric $27$-qubit sixth-order anticoherent state cannot have a
proper rotational symmetry of order $m\ge4$ in its Majorana constellation.

\paragraph{Reduction to moments.}
Rotate the symmetry axis to $z$, which preserves anticoherence. Invariance
of the constellation under a rotation by $2\pi/m$ means that its state ray
is an eigenvector of that rotation. Consequently the Dicke amplitudes
can be supported only on one residue class
\[
 A_{m,a}=\{a,a+m,\ldots\}\cap\{0,\ldots,27\},\qquad0\le a<m.
\]
Write $p_r=|c_r|^2$. Sixth-order anticoherence implies
\[
 \sum_{r=0}^{27}p_r f(r)=\frac1{28}\sum_{r=0}^{27}f(r)
 \qquad(\deg f\le6),
\]
because every polynomial of degree at most six in
$J_z=r-27/2$ decomposes into spin tensors of ranks at most six, whose
non-scalar expectations vanish.

\paragraph{Order four: a closed rational contradiction.}
For $m=4$, every residue class has seven nodes $a+4i$, $i=0,\ldots,6$,
with $a=0,1,2,3$. Its middle-node cardinal polynomial is
\[
 L_a(r)=-\frac{\displaystyle\prod_{h\in\{4,8,12\}}
                   \bigl((r-a-12)^2-h^2\bigr)}{147456}.
\]
It equals one at $r=a+12$ and zero at the other six nodes. Thus
\[
 p_{a+12}=\frac1{28}\sum_{r=0}^{27}L_a(r)=
 \begin{cases}
  -180963/229376,&a=0,3,\\
  -51855/229376,&a=1,2.
 \end{cases}
\]
Every value is negative, contradicting $p_r\ge0$.

\paragraph{Orders five through twenty-seven: exact finite arithmetic.}
For $5\le m\le27$, each nonempty residue class has at most six nodes.
Define its annihilator
\[
 q_{m,a}(r)=\prod_{x\in A_{m,a}}(r-x).
\]
The left-hand moment expectation is zero and $\deg q_{m,a}\le6$.
However,
\[
 \frac1{28}\sum_{r=0}^{27}q_{m,a}(r)\ne0
 \qquad(5\le m\le27,\ 0\le a<m).
\]
The supplied self-contained checker verifies each of these finitely many
rational inequalities with exact integer and rational arithmetic and
records every value in \texttt{symmetry-verification.json}. No numerical
thresholds or floating-point eigenvalues enter this check.
For $m\ge28$, each residue has at most one node. A single Dicke state
would have $\langle r\rangle=r\in\mathbb Z$, whereas first-order
anticoherence requires $\langle r\rangle=27/2$. This is impossible.
Together these arguments prove the claim.

\paragraph{Scope.}
A hypothetical sixth-order state could still have only rotations of
orders two or three, or have no nontrivial rotational symmetry. Nothing
in this argument excludes these possibilities. Failed numerical searches
in those classes are not substituted for a global exclusion proof.
\end{document}
